% Author: Felix Plasser (2025)
%
% This is a simple template for generating a latex supporting information.
% Modify as needed feel free to use for your own research papers.
% If you have improvements, please send as a merge request.

\documentclass[10pt,a4paper]{article}
\usepackage[utf8x]{inputenc}
\usepackage{ucs}
\usepackage{amsmath}
\usepackage{amsfonts}
\usepackage{amssymb}
\usepackage{physics}
\usepackage[left=2cm,right=2cm,top=2cm,bottom=2cm]{geometry}

\renewcommand{\thefigure}{S\arabic{figure}}
\renewcommand{\thetable}{S\arabic{table}}
\renewcommand{\theequation}{S\arabic{equation}}
\renewcommand{\thesection}{}
\renewcommand{\thesubsection}{}

\newcommand{\sititle}[1]{\title{Supporting Information\\~\\\emph{for}\\~\\ \textbf{#1}}}

\begin{document}

\sititle{Very important paper}
\author{M. Mustermann, C. Musterfrau}

\maketitle

\vspace{2em}
\tableofcontents
\listoftables
\listoffigures

\newpage

\section{S1: Mathematical background}
\subsection{Note S1: Derivation of XXX}
First we have to realise that
%
\begin{equation}
1 + 1 = 2
\end{equation}
%
and this will allow us to find that
%
\begin{equation}
\iint \Psi(x,y)\sin(x)\mathrm{d}y \mathrm{d}x = -\frac{\pi}{2}
\end{equation}
%
which then means that
%
\begin{equation}
\bra{\Psi_I}\widehat{\mathrm{H}}\ket{\Psi_J} = \delta_{IJ}
\end{equation}

\subsection{Note S2: Derivation of YYY}

\newpage

\section{S2: Additional results on ...}

\phantom{ } % spacer to make sure things stay in order

\begin{table}[h!]
\centering
\caption{Important results on ....}

\begin{tabular}{cc}
\hline
Header & Header\\
\hline
a & b\\
\hline
\end{tabular}
\end{table}

\phantom{ } % spacer to make sure things stay in order

\begin{table}[h!]
\centering
\caption{Not so important results on ....}

\begin{tabular}{cc}
\hline
Header & Header\\
\hline
c & d\\
\hline
\end{tabular}
\end{table}

\phantom{ } % spacer to make sure things stay in order

\begin{figure}[h!]
\caption{Interesting figure}
%\includegraphics[scale=•]{•}
\end{figure}

\end{document}
